\documentclass[11pt]{article}
\usepackage[margin=1in]{geometry}
\usepackage[T1]{fontenc}
\usepackage{lmodern}
\usepackage{microtype}
\usepackage{xcolor}
\usepackage{hyperref}
\usepackage{booktabs}
\usepackage{enumitem}
\usepackage{listings}
\usepackage{tcolorbox}
\usepackage{parskip}
\usepackage{array}
\newcolumntype{L}[1]{>{\raggedright\arraybackslash}p{#1}}

% \path{} (from url, loaded by hyperref) prints paths and identifiers
% verbatim -- so `--` is not turned into a dash -- and may break at / _ - .
\makeatletter
\g@addto@macro\UrlBreaks{\do\_\do\-\do\.}
\makeatother
\urlstyle{tt}
\setlength{\emergencystretch}{3em}

\hypersetup{
  colorlinks=true,
  linkcolor=blue!60!black,
  urlcolor=blue!60!black
}

\definecolor{codebg}{RGB}{245,245,245}
\definecolor{codecomment}{RGB}{90,100,90}
\definecolor{codekeyword}{RGB}{0,70,140}
\definecolor{codestring}{RGB}{150,60,20}

\lstdefinestyle{pythonstyle}{
  language=Python,
  backgroundcolor=\color{codebg},
  basicstyle=\ttfamily\small,
  keywordstyle=\color{codekeyword}\bfseries,
  commentstyle=\color{codecomment}\itshape,
  stringstyle=\color{codestring},
  showstringspaces=false,
  breaklines=true,
  frame=single,
  rulecolor=\color{black!20},
  columns=fullflexible,
  keepspaces=true
}
\lstdefinestyle{shellstyle}{
  language=bash,
  backgroundcolor=\color{codebg},
  basicstyle=\ttfamily\small,
  showstringspaces=false,
  breaklines=true,
  frame=single,
  rulecolor=\color{black!20},
  columns=fullflexible,
  keepspaces=true
}
\lstset{style=pythonstyle}

\newtcolorbox{keyidea}{
  colback=blue!5, colframe=blue!50!black,
  title=Key Idea, fonttitle=\bfseries
}
\newtcolorbox{tipbox}{
  colback=green!5, colframe=green!45!black,
  title=Tip, fonttitle=\bfseries
}
\newtcolorbox{warningbox}{
  colback=orange!7, colframe=orange!60!black,
  title=Common Pitfall, fonttitle=\bfseries
}

\title{
  \textbf{Codebase Change Summary --- Week 8}\\[4pt]
  \large Functional Testing\\[2pt]
  \normalsize SE 413/513 --- Software Testing
}
\author{Demo Testing Suite}
\date{}

\begin{document}
\maketitle

\begin{center}
\textit{What changed in \path{demo-testing-suite/} since Week 7, and why}
\end{center}

\vspace{0.3cm}

% ============================================================
\section{Purpose of This Document}
% ============================================================
The \path{demo-testing-suite/} codebase is a single, shared project that
grows one lecture at a time rather than being rebuilt from scratch each
week. This document summarizes exactly what was added or changed for
Week 8 and, more importantly, explains \emph{why} each change exists ---
which concept from the Week 8 reading, \textit{Functional Testing:
Designing Tests from the Specification}, it demonstrates. Read this
alongside the reading and \path{demo-testing-suite/README.md}.

\begin{keyidea}
Nothing from Weeks 1--7 was removed or changed. This week adds one
small, self-contained project that runs every example in the reading:
the decision table for shipping fees, the \texttt{Order} state machine,
the pairwise checkout rows, and the Hypothesis property for splitting a
bill --- together with the buggy versions each technique catches. The
main suite is untouched, so its expected results are the same as last
week.
\end{keyidea}

% ============================================================
\section{What Changed At a Glance}
% ============================================================
\begin{center}
\small
\begin{tabular}{@{}L{5.9cm}L{1.5cm}L{7.0cm}@{}}
\toprule
\textbf{File} & \textbf{Status} & \textbf{Purpose} \\
\midrule
\multicolumn{3}{@{}l}{\textit{The reading, runnable (not in the default run): \path{demos/functional_testing/}}}\\
\path{shop/shipping.py}, \path{shop/shipping_buggy.py} & New &
\texttt{shipping\_fee} for REQ-1--4, and a version that checks the
free-shipping rule first. \\
\path{shop/order.py} & New & The \texttt{Order} state machine for
REQ-5--7. \\
\path{shop/split.py}, \path{shop/split_fixed.py} & New &
\texttt{split\_bill} for REQ-8: the reading's buggy first version, and
the fix in integer cents. \\
\path{tests/test_shipping_dt.py} & New & One test per decision-table
rule (R1--R5). \\
\path{tests/test_shipping_bva.py} & New & The \$50 boundary: 49.99
and 50.0. \\
\path{tests/test_order.py} & New & All-transitions coverage: 5 valid,
15 invalid, 2 refund tests. \\
\path{pairwise.py}, \path{tests/test_checkout_pairwise.py} & New & The
\texttt{allpairspy} rows, a check that every pair is covered, and the
rows as \texttt{pytest} cases. \\
\path{tests/test_split_pbt.py} & New & The Hypothesis property for
REQ-8. \\
\texttt{README.md} & New & Every command, in the reading's order, with
its expected result. \\
\midrule
\multicolumn{3}{@{}l}{\textit{Tooling and documentation}}\\
\texttt{requirements.txt} & Modified & Adds \texttt{hypothesis} and
\texttt{allpairspy}. \\
\texttt{.gitignore} & Modified & Ignores \texttt{.hypothesis/},
Hypothesis's local example database. \\
\path{docs/DEMO_GUIDE_week08_functional_testing.md} & New &
Section-by-section lecture companion for the reading. \\
\texttt{README.md} (suite and repository) & Modified & Updated layout
and the Weekly Updates index. \\
\bottomrule
\end{tabular}
\end{center}

\begin{tipbox}
Reinstall the requirements before running anything this week:
\begin{lstlisting}[style=shellstyle]
cd demo-testing-suite
pip install -r requirements.txt
\end{lstlisting}
\end{tipbox}

% ============================================================
\section{The Reading's Examples, Runnable}
% ============================================================
\textbf{Location:} \path{demos/functional_testing/}

This folder is laid out exactly like the reading
(\path{shop/shipping.py}, \path{tests/test_order.py}, \ldots), so the
reading's \texttt{from shop.}\textit{module}\texttt{ import} lines and \texttt{pytest}
commands work as printed. \textbf{Run its commands from inside the
folder:}
\begin{lstlisting}[style=shellstyle]
cd demos/functional_testing
\end{lstlisting}

\begin{center}
\footnotesize
\begin{tabular}{@{}L{5.6cm}L{4.4cm}L{4.6cm}@{}}
\toprule
\textbf{Run} & \textbf{Result} & \textbf{Reading section} \\
\midrule
\texttt{pytest tests/test\_shipping\_bva.py} & 2 passed & Recap: EP and
BVA \\
\ldots{} after changing \texttt{>=~50} to \texttt{>~50} & fails at
\texttt{50.0} & Recap: EP and BVA \\
\texttt{pytest tests/test\_shipping\_dt.py -v} & 5 passed, IDs R1--R5 &
One Test per Rule \\
\ldots{} importing \path{shop.shipping_buggy} & R1 fails; 1 failed, 4
passed & What It Catches \\
\texttt{pytest tests/test\_order.py} & 22 passed & Implementing the
Tests \\
\texttt{python pairwise.py} & 10 rows; 45 of 45 pairs & The Pairwise
Idea \\
\texttt{pytest tests/test\_checkout\_pairwise.py} & 10 skipped (no
oracle yet) & Limitations \\
\texttt{pytest tests/test\_split\_pbt.py} & fails; shrunk to
\texttt{total\_cents=1, people=2} & From Examples to Properties \\
\ldots{} importing \path{shop.split_fixed} & 1 passed & Fixing the
Defect \\
\bottomrule
\end{tabular}
\end{center}

\begin{warningbox}
Two demos are run by changing one import line, and they go in opposite
directions. \path{tests/test_shipping_dt.py} imports the
\emph{correct} \path{shop.shipping}; you switch it to
\path{shop.shipping_buggy}. \path{tests/test_split_pbt.py} imports the
\emph{buggy} \path{shop.split}; you switch it to \path{shop.split_fixed}.
That is why a plain \texttt{pytest} in this folder reports
\texttt{1 failed, 29 passed, 10 skipped}: the one failure is REQ-8's, on
purpose. Switch any import you changed back afterwards.
\end{warningbox}

\subsection*{Why the pairwise test is skipped}
The reading leaves the body of \texttt{test\_checkout\_pairwise} as an
ellipsis (\texttt{...}), a placeholder. Run that way, all ten rows would \emph{pass} while
checking nothing --- Week 7's ``100\% coverage, zero assertions'' trap
in a new form. Pairwise testing chooses combinations; it does not supply
expected results. Until there is a \texttt{checkout()} function and an
expected result for each row, every row is skipped with the reason
``no oracle yet''.

\subsection*{\texttt{.hypothesis/}}
Hypothesis saves every failing input in a local database
(\path{.hypothesis/}) and tries those inputs first on the next run, so
a bug that comes back is caught at once. Delete the folder to watch
Hypothesis search from scratch. It is listed in \texttt{.gitignore}.

% ============================================================
\section{Something to Check Yourself}
% ============================================================
The pairwise slide's ``Done early?'' asks whether every valid pair is
still covered after filtering out ``gift card with a non-USD currency''.
Run the same coverage check that \path{pairwise.py} uses on the filtered
set before you answer.

\begin{keyidea}
A generated test set is a claim, not a guarantee. Last week, a coverage
number needed reading before it could be trusted; this week, the
output of a pairwise generator does. A few lines of checking code ---
like the one in \path{pairwise.py} --- tell you whether the claim holds.
\end{keyidea}

% ============================================================
\section{Updated: Expected Test Results}
% ============================================================
From the project root:
\begin{lstlisting}[style=shellstyle]
cd demo-testing-suite
pip install -r requirements.txt
pytest -v
\end{lstlisting}
you should still see \textbf{72 passed, 2 xfailed, 1 failed}, exactly
as in Week 7. \path{demos/functional_testing/} is not part of that run
(\path{pytest.ini} sets \texttt{testpaths = tests}); run it from inside
its folder.

% ============================================================
\section{How This Connects to the Reading}
% ============================================================
\begin{keyidea}
Weeks 5--7 asked whether the tests exercised the code. This week asks
whether the code does what the specification says. Each file in the
demo folder is one model from the reading --- a table, a state machine,
a parameter set, a property --- turned into tests, and each buggy
version is a defect that the model finds and a coverage report would
not.
\end{keyidea}

% ============================================================
\section{Suggested Next Steps}
% ============================================================
\begin{enumerate}
\item Switch \path{tests/test_shipping_dt.py} to the buggy version and
add a case for rule R2 at \$80. Predict the result before running it.
Then switch the import back.
\item Add the \texttt{return\_item} event to \path{shop/order.py}
(Exercise 3) and extend \path{tests/test_order.py} to keep
all-transitions coverage. How many tests pass now?
\item Add \texttt{currency} to \path{pairwise.py} (Exercise 4), filter
out gift card with a non-USD currency, and run the coverage check on
the filtered rows.
\item Write the two \texttt{apply\_coupon} properties from Exercise 5 in
\path{demos/coverage_reports/}. If Hypothesis finds a counterexample,
decide whether the defect is in the code or in the property.
\item Run
\texttt{pytest -{}-cov=shop -{}-cov-branch -{}-cov-report=term-missing}
on just the two shipping test files, and explain why four
\path{shop/} modules report 0\%.
\end{enumerate}

\end{document}
